\documentclass{article}
\title{個人財務の設計 — お金の数学で暮らしを整える}
\author{TeX 図書館}

\begin{document}

\section{この教科書のために — 生活にファイナンスを接続する}

『ファイナンス基礎』の道具(複利・現在価値・実質リターン)を\textbf{自分の暮らし}に接続するのがこの教科書です。投資手法の前に立つ\textbf{家計の土台}——支出の可視化・緊急資金・インフレ対策・保険の選び方・ライフプランの数理を扱います。

\begin{quote}
\textbf{【心構え】}\ 個人の状況は千差万別で、この本は一般的な数理を示すものであり個別の助言ではありません。判断の最終責任は各自にあります。
\end{quote}

\section{支出の可視化 — 測れないものは管理できない}

家計の第 0 步は\textbf{収入と支出の構造}を記録することです。家計を「固定費」と「変動費」に分けると、数学的な扱いが変わります:

\begin{itemize}
\item \textbf{固定費}(住居・通信・サブスク): 毎月必ず発生し、\textbf{一度削ると毎月効き続ける}。削減効果は複利で効く「元本の削減」に相当
\item \textbf{変動費}(食費・交際費): 日々の選択で動く。削減は「日々の利息」に相当し、意志力への依存度が高い
\end{itemize}

\begin{quote}
\textbf{【例】}\ 月 1000 円の使っていないサブスクを解約する効果。年 12000 円、年利 5\% で 20 年間投資に回すと \(12000 \times \frac{1.05^{20} - 1}{0.05} \approx 39.7\) 万円。\textbf{月 1000 円の判断が 20 年後に 40 万円}——固定費の削減は金額以上に重い。
\end{quote}

\begin{quote}
\textbf{【演習】}\ 月 3000 円の節約を年利 5\% で 10 年間積み立てた場合の将来価値を求めよ(年金の将来価値の公式)。
\end{quote}

\textbf{解答の指針:}\ 年 36000 円 × \(\frac{1.05^{10} - 1}{0.05} \approx 12.58\) → 約 45.3 万円。

\section{貯蓄率 — 財務自由への単一の最重要変数}

\textbf{貯蓄率} = (収入 − 支出) ÷ 収入。これが家計の数理で最も強力なレバーです。

\begin{quote}
\textbf{【定理:貯蓄率と労働年数】}\ 「資産から年支出の 4\% を取り崩して生活できる状態」(\textbf{4\% ルール})を目標にすると、必要資産は\textbf{年支出の 25 倍}。貯蓄率 \(s\) のときの労働年数は概算で
\[ n \approx \frac{\ln\left(1 + \frac{0.04 \times 25(1-s)}{s}\right)}{\ln(1 + r)} \approx \frac{\ln\left(1 + \frac{1-s}{s}\right)}{\ln(1+r)} \quad (r: \text{実質リターン}) \]
の形になり、\textbf{貯蓄率が高いほど年数は急激に短くなる}(非線形!)。
\end{quote}

\begin{center}
\begin{tabular}{|c|c|}
\hline
貯蓄率 & 労働年数の目安(実質 4\%) \\
\hline
10\% & 約 51 年 \\
25\% & 約 32 年 \\
40\% & 約 22 年 \\
60\% & 約 12.5 年 \\
\hline
\end{tabular}
\end{center}

\begin{quote}
\textbf{【演習】}\ 貯蓄率 10\% と 60\% で労働年数が 4 倍以上変わる構造を、「支出が目標額と貯蓄速度の両方を決める」観点から説明せよ。
\end{quote}

\textbf{解答の指針:}\ 貯蓄を増やすことは\textbf{必要額(年支出の 25 倍)を減らす}と\textbf{積立速度を上げる}の両方に同時に効く二重効果があるため、効果が非線形に現れます。

\section{緊急資金 — 悪事の数学に対する防壁}

\textbf{緊急資金}(生活費の 3〜6 か月分を現金・普通預金で保有)は、投資の前段階の必須装備です。数理的根拠:

\begin{enumerate}
\item \textbf{破産確率の概念}(『投資の数学』第 7 節)の家計版——突然の失業・疾病は「ゼロに吸収される」状態で、回復が最も困難
\item 緊急時に\textbf{投資資産を下値で売る}ことを避ける(一時損失の確定 + 複利エンジンの部品損失)
\item 現金の機会損失(実質リターン負)は「保険料」として払う合理性がある
\end{enumerate}

\begin{quote}
\textbf{【演習】}\ 月支出 20 万円の人が 4 か月分(80 万円)の緊急資金を持つ意義を、「低頻度・大損害」の対処という観点から保険との共通点で述べよ。
\end{quote}

\textbf{解答の指針:}\ どちらも\textbf{大きな損失イベントに備えた事前の準備}。頻度は低いが発生時の打撃が構造的なものに対して、確率の低さを理由に無防備でいるのが最大のリスクです。

\section{インフレ — 現金への静かな課金}

インフレ率 \(\pi\) のとき、現金の実質リターンは約 \(-\pi\)。\textbf{72 の法則}によれば、インフレ 3\% なら 24 年で現金の購買力は半減します。

\begin{center}
\begin{tabular}{|c|c|}
\hline
インフレ率 & 購買力が半減する年数 \\
\hline
1\% & 約 70 年 \\
2\% & 約 36 年 \\
3\% & 約 24 年 \\
5\% & 約 14 年 \\
\hline
\end{tabular}
\end{center}

\begin{quote}
\textbf{【演習】}\ 「安全のために全額現金」戦略の 30 年後を、インフレ 2\%・運用 0\% で評価せよ。
\end{quote}

\textbf{解答の指針:}\ 購買力は \(1.02^{-30} \approx 0.55\) 倍。\textbf{元本は減っていないのに実質 45\% 減}。「減らないこと」は「価値を保つこと」ではない。

\section{保険 — 負の期待値を払う合理性}

保険は数学的には\textbf{期待値が負の取引}です(保険料に運営費・利益が上乗せされるため)。ではなぜ合理的か——\textbf{破産確率を消すため}です。

\begin{quote}
\textbf{【定理:保険の合理条件】}\ 頻度が低く・損害が大きく・生活に致命的なリスク(火災・死亡・医療の大型手術)には保険が合理的。頻度が高く・損害が小さいリスク(スマホ破損など)への保険は、負の期待値を高頻度で払うことになるため非合理的。
\end{quote}

\begin{quote}
\textbf{【演習】}\ 「年間保険料 3 万円で 3000 万円保障」の死亡保険と「年 5000 円で 5 万円保障」の携帯保険。それぞれの 1 円あたりの期待構造を、負の期待値の大きさの観点で比較せよ。
\end{quote}

\textbf{解答の指針:}\ 死亡保険は「極低頻度・超大損害」を薄く広くカバーする(負の期待値幅は小さい)。携帯保険は「高頻度・小損害」に重いコストが乗る構造。損害が自己資金で吸収できるなら保険を付けないのが数理的には正。

\section{ライフプラン — 現在価値で人生を見る}

人生の大きな支出・収入を\textbf{現在価値}に揃えると、意思決定が明快になります。

\begin{enumerate}
\item \textbf{住居}: 賃貸 vs 購入は「現在価値の比較」問題(家賃支払いの PV vs 購入費・維持費・売却時価の PV)。金利・居住期間・資産価値の変動が結論を左右する
\item \textbf{教育投資}: 自己投資のリターン = 賃上げ額 × 勤続年数の PV − 費用の PV。回収期間を数値で見積もる
\item \textbf{退職資金}: 必要額 = 退職後年支出 × 25(4\% ルール)に公的年金の現価を差し引いた額
\end{enumerate}

\begin{quote}
\textbf{【演習】}\ 退職後の年支出 300 万円、公的年金の現価 1500 万円とする。4\% ルールで必要な自主運用資産を求めよ。
\end{quote}

\textbf{解答の指針:}\ 300 × 25 = 7500 万円 − 1500 万円 = 6000 万円。\textbf{目標額を現在価値の言葉で分解}すると、各要素の寄与と達成可能性が見える。

\section{行動の設計 — 数学を実行に移す仕組み}

最後の、そして最難関のパートは\textbf{実行}です。家計の数理は単純でも、人間は単純に動きません。仕組みで解決します:

\begin{enumerate}
\item \textbf{先取り自動化}: 給与天引き・自動積立で「残りから貯める」ではなく「貯めた残りで暮らす」。意志力を使わない設計
\item \textbf{ルールの文書化}: 配分(たとえば緊急資金・インデックス・個別)の割合と、増額・売却の条件を事前に書いておく
\item \textbf{定期点検}: 年 1 回、収支・資産配分・保険の見直し。日々の価格を見ない(見ても行動しない)のがポイント
\item \textbf{情報の節制}: 日次の価格ニュースは判断精度を上げず不安だけを増やします(『投資の数学』の行動バイアス節の再登場)
\end{enumerate}

\begin{quote}
\textbf{【演習】}\ 「市場が 20\% 下落したら積立額を倍にする」ルールを数理的に評価せよ(ドルコスト平均法の拡張として)。
\end{quote}

\textbf{解答の指針:}\ 下落時に安値で多く買う構造で、平均購入単価を下げる効果が理論上ある。ただし下落が続く間の現金枯渇と、下落判断(20\% の定義)を機械化できるかが実務の鍵。

\section{おわりに — 数字と暮らしの間}

この教科書の骨格: 支出の可視化(§2)、貯蓄率の非線形パワー(§3)、緊急資金(§4)、インフレ対策(§5)、保険の合理条件(§6)、現在価値で見るライフプラン(§7)、実行の仕組み化(§8)。

『ファイナンス基礎』の道具が生活の意思決定に、『投資の数学』のリスク管理が資産運用に接続され、\textbf{生活設計の全体}が数理で語れるようになりました。数字は暮らしを冷たくするものではなく、\textbf{将来の自分と交渉するための言語}です。

\begin{quote}
\textbf{【総合演習】}
\begin{enumerate}
\item 月 2 万円を年利 4\% で 15 年積立した将来価値を求めよ。
\item インフレ 2\% のもとで実質リターン 3\% を得るには名目リターン何\% が必要か。
\item 貯蓄率 50\% の人が 4\% ルールで財務自由に至る概算年数を述べよ。
\end{enumerate}
\end{quote}

\textbf{解答の指針:}\ 1. 年 24 万円 × \(\frac{1.04^{15} - 1}{0.04} \approx 20.15\) → 約 484 万円。2. \((1.03)(1.02) - 1 \approx 5.06\%\)。3. 約 17 年程度(表の 60\% の 12.5 年と 40\% の 22 年の間)。

\end{document}
